\section{The classification problem}

For $s\in\{1/6,1/4,1/3,1/2\}$ let
\[
 F_s(z)={}_3F_2\!\left(\begin{matrix}1/2,s,1-s\\1,1\end{matrix};z\right)
       =\sum_{n\ge0}c_s(n)z^n,
 \qquad c_s(n)=\frac{(1/2)_n(s)_n(1-s)_n}{(n!)^3},
 \qquad\thetaop=z\frac d{dz}.
\]
The pairs $(s_\ell,R_\ell)$ attached to levels $\ell=1,2,3,4$ are
\[
 (1/6,1728),\qquad(1/4,256),\qquad(1/3,108),\qquad(1/2,64).
\]
We use the power-series branch at zero and define
\begin{equation}\label{eq:series}
 \mathfrak S_\ell(A,B,C)
  =\sum_{n\ge0}R_\ell^nc_{s_\ell}(n)(A+Bn)C^{-n}
  =(A+B\thetaop)F_{s_\ell}(R_\ell/C),
\end{equation}
where the last expression always means the ordinarily convergent sum.

\begin{definition}
A primitive standard identity is an ordinarily convergent equality
\[
 \mathfrak S_\ell(A,B,C)=\frac\alpha\pi,
 \qquad A,B,C\in\Z,\quad C\ne0,\quad\gcd(A,B)=1,
 \quad\alpha\in\Qbar.
\]
Classes identify $(A,B,\alpha)$ with $(-A,-B,-\alpha)$, at the same
level and denominator. The multiplier is initially allowed to be zero.
\end{definition}

The primitive condition concerns the coefficient pair. Requiring only
$\gcd(A,B,C)=1$ would permit infinitely many multiples of the identity
$(A,B,C)=(1,6,256)$ at level four. Analytic continuation is likewise
excluded from the definition: at a convergence boundary, the value must
be the sum of the original series.

\begin{hypothesis}[Quadratic-period independence]\label{hyp:QPI}
Let $z=R_\ell/C\in[-1,1)\setminus\{0\}$, with $C\in\Z$, and put
$x=(1-\sqrt{1-z})/2$. For the algebraic elliptic curve $E_{s_\ell,x}$
of Section~\ref{sec:families}, use the analytic cycle normalized by
\[
 \omega=\frac{2\pi}{\sqrt6}
 {}_2F_1(s_\ell,1-s_\ell;1;x),
 \qquad h=\int_\gamma X\,\frac{dX}{y}.
\]
If this curve is non-CM, then $\pi,\omega^2,\omega h$ are linearly
independent over $\Qbar$.
\end{hypothesis}

\begin{theorem}[Conditional classification]\label{thm:main-conditional}
Assume Hypothesis~\ref{hyp:QPI}. Every primitive standard identity has
a CM elliptic fiber. There are exactly $36$ classes, comprising
$11,12,9,4$ classes at levels $1,2,3,4$, respectively, as listed in
Tables~\ref{tab:standard1}--\ref{tab:standard4}.
For each fixed level and denominator, the primitive coefficient pair
is unique up to simultaneous sign, and its multiplier is nonzero.

For the linear Euler companion $F_s(z)/(1-z)$, the complete list has
$35$ classes. It is obtained by the coefficient transport and ordinary
convergence criterion of Theorem~\ref{thm:euler}.
\end{theorem}
\begin{proof}
The elliptic realization and its first derivative identify the
coefficient space with de Rham cohomology. In the normalized basis,
\begin{equation}\label{rev:eq:leading-reduction}
 \pi\mathfrak S_\ell(A,B,C)
  =\frac{3}{2\pi}\bigl((A+Bu_s)\omega^2+Bv_s\omega h\bigr),
 \qquad u_s,v_s\in\Qbar,\quad v_s\ne0.
\end{equation}
Thus a nonzero coefficient pair gives a nontrivial quadratic-period
relation. Hypothesis~\ref{hyp:QPI} forces CM; the positive endpoint is
excluded unconditionally in Section~\ref{sec:endpoint}. CM reciprocity
and the finite modular correspondence then bound the class number and
give the tabulated parameters. The complementary CM eigenline supplies
the identity, and its equivariance supplies a rational coefficient
slope. These statements, with the exact finite enumeration, are proved
in Section~\ref{sec:enumeration}. Chudnovsky's theorem gives uniqueness
and nonvanishing in Section~\ref{sec:unique}. Finally, the invertible
Euler transport gives the companion classification.
\end{proof}

\begin{theorem}[Unconditional uniqueness]\label{thm:unconditional}
At every nonzero algebraic real $z\in[-1,1)$, there is at most one
$[A:B]\in\mathbf P^1(\Qbar)$ such that
$(A+B\thetaop)F_s(z)\in\Qbar/\pi$. A nonzero coefficient pair cannot
give zero. At $z=1$, no nonzero algebraic pair gives a convergent
identity in $\Qbar/\pi$.
\end{theorem}

The first two assertions are consequences of the identity-line theorem;
the endpoint uses CM period independence. In particular, the uniqueness
assertion in the leading theorem needs no conjecture.

The main proofs are organized by the information carried by an identity.
Section~\ref{sec:general-period-square} constructs its coefficient line
intrinsically. Section~\ref{sec:hinge} describes the precise extra
comparison information that would force CM and proves the implications
from named period conjectures. Section~\ref{sec:enumeration} completes
the arithmetic classification within the CM locus. The arithmetic
criteria and scarcity theorems then give unconditional conclusions with
different hypotheses: they neither assume nor prove the missing
single-value implication. The appendices contain coordinate
certificates, refinements, further families, and obstructions to
strengthening those hypotheses.
